% =========================================================================
\section{La ley de Coulomb}\label{sec:coulomb}
% =========================================================================
% -------------------------------------------------------------------------
\subsection{La carga eléctrica}\label{sec:carga}
% -------------------------------------------------------------------------
La materia está formada por cargas eléctricas positivas y negativas que, en la materia
ordinaria, se compensan: los cuerpos son neutros. Recoge este equilibrio el siguiente
ejemplo, que da una idea de la enorme intensidad de la fuerza eléctrica.

\begin{ejemplo}[¿Cuánta carga sobra si se rompe la neutralidad?]\label{ej:neutralidad}
Una persona de $70\,\mathrm{kg}$ contiene unos $2\cdot10^{28}$ electrones. Si cada una de
dos personas situadas a $1\,\mathrm{m}$ tuviese un 1\,\% más de electrones que de protones,
cada una llevaría una carga $Q\approx2\cdot10^{26}\times1{,}6\cdot10^{-19}\,\mathrm{C}\approx3\cdot10^{7}\,\mathrm{C}$,
y la fuerza de repulsión sería
\[
  F\approx\frac{1}{4\pi\eo}\frac{Q^{2}}{d^{2}}\approx9\cdot10^{9}\cdot\frac{(3\cdot10^{7})^{2}}{1^{2}}\,\mathrm{N}\sim10^{25}\,\mathrm{N},
\]
del orden del peso de toda la Tierra ($\sim6\cdot10^{25}\,\mathrm{N}$).
\end{ejemplo}

Se puede suministrar o extraer carga a un cuerpo neutro, por ejemplo frotando unos materiales
con otros (vidrio con seda): el vidrio queda cargado positivamente (con defecto de electrones) y la
seda negativamente (con exceso). Pero la carga total del sistema se \emph{conserva}: la carga no se
crea ni se destruye. Además, está \emph{cuantizada}: cualquier carga $q$ es un múltiplo entero
de la carga elemental, $e=1{,}6\cdot10^{-19}\,\mathrm{C}$ (carga del electrón, $-e$). Entre dos cuerpos cargados aparecen
fuerzas que dependen de la cantidad de carga y de cómo está distribuida sobre el volumen o la superficie
de cada cuerpo. Por eso es útil la idealización de \emph{carga puntual}: toda la carga concentrada
en un punto geométrico. Es una buena aproximación cuando la distancia entre los cuerpos es mucho mayor que
sus dimensiones.

% -------------------------------------------------------------------------
\subsection{Evidencia experimental: de la balanza de torsión a Cavendish}\label{sec:coulomb-exp}
% -------------------------------------------------------------------------
La primera investigación cuantitativa de la dependencia de esta fuerza con la distancia fue
realizada por Coulomb en 1785~\cite{COULOMB} con una \emph{balanza de torsión}
(\cref{fig:torsion}). Una aguja aislante, colgada de un hilo fino, lleva una esfera cargada en un extremo
y un contrapeso en el otro; otra esfera cargada, fija, repele a la primera. En el equilibrio, el par
de torsión del hilo, proporcional al ángulo girado, iguala al par de la fuerza electrostática: midiendo ángulos se miden
fuerzas. Es un experimento muy delicado, sujeto a pérdidas de carga y a fuerzas de inducción sobre las esferas, y de él
solo se obtiene con precisión modesta la ley $1/r^{2}$.

% Figura: balanza de torsión de Coulomb (esquema).
\begin{figure}[ht]
\centering
\tikzset{esf/.style={circle, inner sep=0pt, minimum size=9pt, ball color=azulacero!75, draw=azulprof, line width=0.4pt},
         esffija/.style={circle, inner sep=0pt, minimum size=9pt, ball color=acento!80, draw=acento!60!black, line width=0.4pt}}
\begin{minipage}[t]{0.56\textwidth}\centering
\begin{tikzpicture}[baseline=(current bounding box.north), x=1cm, y=1cm]
  % (a) vista en perspectiva: cilindro de vidrio, hilo, aguja y esferas
  \def\rx{1.55}\def\ry{0.42}
  \fill[azulclaro!30] (-\rx,0) -- (-\rx,3.2) arc (180:360:{\rx} and \ry) -- (\rx,0) arc (0:-180:{\rx} and \ry) -- cycle;
  \draw[azulprof, line width=0.8pt] (-\rx,0) -- (-\rx,3.2) (\rx,0) -- (\rx,3.2);
  \draw[azulprof, line width=0.8pt, fill=azulclaro!45] (0,3.2) ellipse ({\rx} and \ry);
  \draw[azulprof, line width=0.8pt] (-\rx,0) arc (180:360:{\rx} and \ry);
  \draw[azulprof, line width=0.5pt, densely dashed] (-\rx,0) arc (180:0:{\rx} and \ry);
  % escala graduada alrededor del cilindro (marcas en la mitad frontal)
  \draw[azulprof!70, line width=0.5pt] (-\rx,1.1) arc (180:360:{\rx} and \ry);
  \foreach \a in {190,200,...,350}{\draw[azulprof!70, line width=0.35pt] ({\rx*cos(\a)},{1.1+\ry*sin(\a)}) -- ({\rx*cos(\a)},{1.1+\ry*sin(\a)+0.1});}
  % cabezal micrométrico
  \draw[azulprof, line width=0.7pt, fill=azulclaro!70] (-0.45,3.52) rectangle (0.45,3.95);
  \draw[azulprof, line width=0.7pt, fill=azulclaro!40] (0,3.95) ellipse (0.45 and 0.13);
  \draw[acento, line width=0.9pt] (0,3.95) -- (0.3,3.99);
  \draw[azulprof, line width=0.7pt, fill=azulclaro!55] (0,3.52) ellipse (0.45 and 0.13);
  % hilo de torsión
  \draw[grispizarra, line width=0.6pt] (0,3.38) -- (0,1.45);
  % aguja (en perspectiva) con esfera móvil y contrapeso
  \draw[grispizarra, line width=1.3pt] (-1.15,1.28) -- (1.15,1.62);
  \fill[azulprof] (0,1.45) circle (1.3pt);
  \node[esf] at (1.15,1.62) {}; \node[esf] at (-1.15,1.28) {};
  % esfera fija sujeta desde la tapa
  \draw[grispizarra, line width=0.9pt] (1.02,3.05) -- (1.02,1.35);
  \node[esffija, minimum size=8pt] at (1.02,1.30) {};
  % rótulos
  \draw[grispizarra, line width=0.3pt] (0.5,3.8) -- (2.3,3.8); \node[font=\footnotesize, anchor=west] at (2.35,3.8) {cabezal micrométrico};
  \draw[grispizarra, line width=0.3pt] (0.02,2.6) -- (2.3,2.6); \node[font=\footnotesize, anchor=west] at (2.35,2.6) {hilo de torsión};
  \draw[grispizarra, line width=0.3pt] (1.25,1.7) -- (2.3,1.95); \node[font=\footnotesize, anchor=west] at (2.35,1.95) {esfera móvil};
  \draw[grispizarra, line width=0.3pt] (1.12,1.18) -- (2.3,1.05); \node[font=\footnotesize, anchor=west] at (2.35,1.05) {esfera fija};
  \draw[grispizarra, line width=0.3pt] (1.5,0.62) -- (2.3,0.35); \node[font=\footnotesize, anchor=west] at (2.35,0.3) {escala graduada};
  \node[font=\footnotesize, anchor=east] at (-1.75,1.1) {aguja};
  \draw[grispizarra, line width=0.3pt] (-1.7,1.15) -- (-1.25,1.27);
\end{tikzpicture}\\[2pt]{\small\scshape (a)}
\end{minipage}\hfill
\begin{minipage}[t]{0.40\textwidth}\centering
\begin{tikzpicture}[baseline=(current bounding box.north), x=1cm, y=1cm]
  % (b) vista superior: escala graduada y ángulo de torsión
  \def\Rg{1.65}
  \draw[azulprof, line width=0.8pt] (0,0) circle (\Rg);
  \foreach \a in {0,6,...,354}{\draw[azulprof!70, line width=0.3pt] (\a:\Rg) -- (\a:{\Rg-0.09});}
  \foreach \a in {0,30,...,330}{\draw[azulprof, line width=0.5pt] (\a:\Rg) -- (\a:{\Rg-0.17});}
  \draw[verdeteo, densely dashed, line width=0.5pt] (0,0) -- (0:1.3);
  \draw[grispizarra, line width=1.3pt] (200:1.2) -- (22:1.2);
  \fill[azulprof] (0,0) circle (1.4pt) node[below left, font=\scriptsize] {$O$};
  \node[esf] at (22:1.2) {}; \node[esf] at (200:1.2) {};
  \node[esffija] at (0:1.3) {};
  \draw[verdeteo, line width=0.6pt] (0:0.7) arc (0:22:0.7);
  \node[font=\small, text=verdeteo] at (11:0.95) {$\theta$};
  \node[font=\scriptsize, text=acento, anchor=west] at (1.85,-0.35) {fija};
  \node[font=\scriptsize, text=azulprof, anchor=south west] at (1.35,1.25) {móvil};
\end{tikzpicture}\\[2pt]{\small\scshape (b)}
\end{minipage}
\caption{Balanza de torsión de Coulomb (esquema). (a)~Una aguja aislante suspendida de un hilo fino lleva una esfera móvil en un extremo y un contrapeso en el otro; una segunda esfera
fija se introduce desde la tapa. Un cabezal micrométrico permite girar el extremo superior del hilo, y una escala graduada rodea el cilindro de vidrio. (b)~Vista desde arriba: al cargar ambas esferas, la repulsión desvía la aguja un
ángulo $\theta$ medido sobre la escala; en el equilibrio, el par recuperador del hilo, proporcional a $\theta$, iguala al de la fuerza electrostática. Girando el cabezal se varía la distancia entre las esferas y se comprueba cómo escala la fuerza con ella.}
\label{fig:torsion}
\end{figure}


Mucho antes que Coulomb, y con mayor precisión, Cavendish había contrastado el exponente
por un camino indirecto (1771--73), que merece la pena entender porque es una de las pruebas más
limpias de la ley (\cref{fig:cavendish}). El argumento descansa en un resultado geométrico: una
cáscara esférica cargada uniformemente ejerce una fuerza nula en su interior \emph{si y solo si} la fuerza
decae como $1/r^{2}$. Este hecho aparece ya, sin demostración, en Franklin y Priestley, y es el análogo
eléctrico del teorema de Newton para la gravedad.

\begin{proposicion}[Argumento de Cavendish]\label{prop:cavendish}
Sea una cáscara esférica con densidad superficial de carga uniforme $\sigma$ y supongamos que
la fuerza entre cargas puntuales es central, repulsiva (para cargas iguales) y de módulo
$K\,q_1q_2/R^{\,n}$. Una carga $q$ en cualquier punto $P$ interior a la cáscara sufre fuerza neta nula
si y solo si $n=2$.
\end{proposicion}
\begin{proof}
Tracemos desde $P$ un cono estrecho de ángulo sólido $\dOm$ y su prolongación opuesta
(\cref{fig:cavendish}b). Los dos conos recortan en la cáscara sendos elementos
$\dif S_1$ y $\dif S_2$ a distancias $r_1$ y $r_2$ de $P$. El ángulo $\alpha$ entre la cuerda y la
normal a la esfera (el radio) es el mismo en ambos extremos, porque el triángulo formado por los dos
radios y la cuerda es isósceles. Por tanto, $\dif S_i=r_i^{2}\,\dOm/\cos\alpha$ y las fuerzas
sobre $q$ actúan a lo largo de la cuerda y en sentidos opuestos,
\[
  \dif F_i=K\,q\,\frac{\sigma\,\dif S_i}{r_i^{\,n}}=\frac{K q\sigma\,\dOm}{\cos\alpha}\;r_i^{\,2-n},
  \qquad
  \dif F_{\mathrm{neta}}=\frac{K q\sigma\,\dOm}{\cos\alpha}\left(r_2^{\,2-n}-r_1^{\,2-n}\right).
\]
Para $n=2$ esta cantidad se anula para cada par de elementos opuestos, y la fuerza total, suma
(integral) sobre todos ellos, es nula. Si $n\neq2$ y $P$ no es el centro, $r_1\neq r_2$ y
$r_2^{\,2-n}\neq r_1^{\,2-n}$, de modo que la contribución de cada pareja no es nula y la
fuerza total es, en general, distinta de cero.
\end{proof}

% Figura: experimento de Cavendish y argumento de la cáscara esférica.
\begin{figure}[ht]
\centering
\tikzset{
  parche/.style={draw=azulprof, line width=0.6pt, fill=azulclaro!70, fill opacity=0.9},
  parche2/.style={draw=acento, line width=0.6pt, fill=acento!25, fill opacity=0.9},
  vn/.style={-{Stealth[length=5pt,width=4.5pt]}, line width=1.1pt}}
\begin{minipage}[t]{0.50\textwidth}\centering
\begin{tikzpicture}[baseline=(current bounding box.north), x=0.82cm, y=0.82cm]
  % (a) cáscara de dos hemisferios con esfera interior unida por un hilo conductor
  \def\Re{1.9}\def\Ri{0.6}
  \shade[ball color=azulclaro!55, opacity=0.35] (0,0) circle (\Re);
  \draw[azulprof, line width=0.8pt] (0,0) circle (\Re);
  \draw[azulprof, line width=0.8pt] (-\Re,0) arc (180:360:{\Re} and 0.45);
  \draw[azulprof, line width=0.5pt, densely dashed] (-\Re,0) arc (180:0:{\Re} and 0.45);
  \shade[ball color=azulacero!70, opacity=0.95] (0,0) circle (\Ri);
  \draw[azulprof, line width=0.7pt] (0,0) circle (\Ri);
  % hilo conductor vertical con collar aislante en la cáscara y terminal exterior
  \draw[grispizarra, line width=1.0pt] (0,\Ri) -- (0,2.55);
  \draw[azulprof, line width=0.6pt, fill=white] (0,\Re) ellipse (0.13 and 0.05);
  \fill[acento] (0,2.55) circle (2.0pt);
  % rótulos con líneas de llamada rectas
  \node[font=\footnotesize, anchor=west] at (0.35,2.55) {hilo conductor};
  \draw[grispizarra, line width=0.3pt] (-2.65,0.45) -- (-1.85,0.45); \node[font=\footnotesize, anchor=east] at (-2.7,0.45) {hemisferios};
  \draw[grispizarra, line width=0.3pt] (2.7,-0.9) -- (1.5,-0.9); \node[font=\footnotesize, anchor=west] at (2.75,-0.9) {cáscara};
  \draw[grispizarra, line width=0.3pt] (2.7,0.0) -- (0.62,0.0); \node[font=\footnotesize, anchor=west] at (2.75,0.0) {esfera interior};
  \draw[verdeteo, <->, line width=0.5pt] (0,0) -- (-35:\Ri) node[pos=0.7, below=1pt, font=\scriptsize, text=verdeteo] {$a$};
\end{tikzpicture}\\[2pt]{\small\scshape (a)}
\end{minipage}\hfill
\begin{minipage}[t]{0.48\textwidth}\centering
\begin{tikzpicture}[baseline=(current bounding box.north), x=1.05cm, y=1.05cm]
  % (b) cono desde un punto interior P que recorta dos elementos opuestos de la cáscara
  \def\Rr{2.0}
  \draw[azulprof, line width=0.9pt, name path=circ, fill=azulfondo, fill opacity=0.5] (0,0) circle (\Rr);
  \coordinate (O) at (0,0); \fill[azulprof] (O) circle (1.2pt); \node[font=\scriptsize, anchor=north east] at (O) {$C$};
  \coordinate (P) at (0.7,-0.25);
  \def\dd{24}\def\ap{9}
  \foreach \dir/\nm in {\dd-\ap/a,\dd+\ap/b,\dd/m}{
    \path[name path=ray\nm] (P) -- ++({\dir}:5);
    \path[name path=rev\nm] (P) -- ++({\dir+180}:5);
    \path[name intersections={of=circ and ray\nm, by=A\nm}];
    \path[name intersections={of=circ and rev\nm, by=B\nm}];}
  \fill[azulclaro, opacity=0.40] (P) -- (Aa) -- (Ab) -- cycle;
  \fill[azulclaro, opacity=0.40] (P) -- (Ba) -- (Bb) -- cycle;
  \draw[azulacero, line width=0.45pt] (Aa) -- (Ba) (Ab) -- (Bb);
  \draw[azulprof, line width=0.8pt] (Bm) -- (Am);
  \draw[acento, line width=2.8pt] (Aa) -- (Ab);
  \draw[acento, line width=2.8pt] (Ba) -- (Bb);
  \draw[verdeteo, line width=0.6pt, densely dashed] (0,0) -- (Am) (0,0) -- (Bm);
  \fill[azulprof] (P) circle (1.5pt) node[below=2pt, font=\small] {$P$};
  \node[font=\scriptsize, text=azulprof, anchor=south east] at ($(P)!0.55!(Am)$) {$r_2$};
  \node[font=\scriptsize, text=azulprof, anchor=north west] at ($(P)!0.55!(Bm)$) {$r_1$};
  \node[font=\small, text=acento, anchor=south west] at ($(Am)+(0.1,0.12)$) {$\dif S_2$};
  \node[font=\small, text=acento, anchor=north east] at ($(Bm)+(-0.12,-0.2)$) {$\dif S_1$};
  % ángulo alpha entre la cuerda y el radio en cada extremo (arcos dibujados a mano)
  \draw[verdeteo, line width=0.55pt] let \p1=(Am), \n1={atan2(\y1,\x1)+180} in
     (Am) ++(\n1:0.45) arc[start angle=\n1, end angle=\dd+180, radius=0.45];
  \draw let \p1=(Am), \n1={atan2(\y1,\x1)+180} in
     node[font=\scriptsize, text=verdeteo] at ($(Am)+({(\n1+\dd+180)/2}:0.78)$) {$\alpha$};
  \draw[verdeteo, line width=0.55pt] let \p1=(Bm), \n1={atan2(\y1,\x1)+180} in
     (Bm) ++(\n1:0.45) arc[start angle=\n1, end angle=\dd, radius=0.45];
  \draw let \p1=(Bm), \n1={atan2(\y1,\x1)+180} in
     node[font=\scriptsize, text=verdeteo] at ($(Bm)+({(\n1+\dd)/2}:0.78)$) {$\alpha$};
\end{tikzpicture}\\[2pt]{\small\scshape (b)}
\end{minipage}
\caption{(a)~Experimento de Cavendish: una esfera conductora interior se conecta, mediante un hilo, a una cáscara formada por dos hemisferios, y se cargan; al retirar el hilo y separar los hemisferios,
la esfera interior resulta descargada. (b)~Argumento geométrico: un cono de ángulo sólido $\dOm$ con vértice en un punto interior $P$ recorta dos elementos, $\dif S_1$ y $\dif S_2$,
a distancias $r_1$ y $r_2$ de $P$; ambos forman con la cuerda el mismo ángulo $\alpha$ (el triángulo formado por los dos radios y la cuerda es isósceles).}
\label{fig:cavendish}
\end{figure}


Cavendish consiguió la configuración de la \cref{fig:cavendish}a: una esfera conductora dentro de
una cáscara, unidas por un hilo, de modo que ambas forman un único conductor. Se cargan, se retira
el hilo y se abre la cáscara. Si la ley fuese exactamente $1/r^{2}$, la carga residiría
toda en la cáscara y la esfera interior quedaría descargada, como en efecto observó: con electrómetros de
la época concluyó que el exponente era $2\pm0{,}02$. Los experimentos posteriores, siguiendo la misma idea, han
reducido la incertidumbre en muchos órdenes de magnitud. Si escribimos el exponente como $2+\varepsilon$:
\begin{center}\small
\renewcommand{\arraystretch}{1.25}
\begin{tabular}{@{}lcc@{}}
\toprule
Experimento & Año & $\lvert\varepsilon\rvert<$\\
\midrule
Cavendish~\cite{CAVENDISH-ELEC} & 1773 & $0{,}02$\\
Maxwell~\cite{MAXWELL} & 1873 & $5\cdot10^{-5}$\\
Plimpton y Lawton~\cite{PLIMPTON} & 1936 & $2\cdot10^{-9}$\\
Williams, Faller y Hill~\cite{WILLIAMS} & 1971 & $\sim3\cdot10^{-16}$\\
\bottomrule
\end{tabular}
\end{center}
El exponente 2 es, así, una de las constantes mejor comprobadas de la física~\cite{FEYNMAN}; sin embargo, la ley de Coulomb
sigue siendo una ley \emph{experimental}, que no admite demostración desde otros principios dentro de
la electrostática.

% -------------------------------------------------------------------------
\subsection{Enunciado de la ley}\label{sec:coulomb-ley}
% -------------------------------------------------------------------------
Elegido un sistema de coordenadas y dos cargas puntuales $q_1$ y $q_2$ en reposo, en el vacío, situadas
en $\vc{r}_1$ y $\vc{r}_2$ (\cref{fig:coulomb}a), la ley de Coulomb afirma que la fuerza es proporcional al
producto de las cargas, inversamente proporcional al cuadrado de la distancia que las separa, está dirigida
a lo largo de la recta que las une, y es atractiva si las cargas son de signos opuestos y repulsiva si
tienen el mismo signo. Matemáticamente:

\begin{definicion}[Ley de Coulomb]\label{def:coulomb}
La fuerza que la carga puntual $q_1$ ejerce sobre la carga puntual $q_2$, en reposo y en el vacío, es
\begin{equation}\label{eq:coulomb}
  \boxed{\;\vc{F}_{12}=K\,\frac{q_1q_2}{R_{12}^{2}}\,\uh{R_{12}}
  =\frac{1}{4\pi\eo}\,\frac{q_1q_2}{\norm{\vc{R}_{12}}^{3}}\,\vc{R}_{12},\qquad
  \vc{R}_{12}=\vc{r}_2-\vc{r}_1\;}
\end{equation}
con $\uh{R_{12}}=\vc{R}_{12}/R_{12}$ y $R_{12}=\norm{\vc{R}_{12}}$. Si $q_1q_2>0$ la fuerza es
repulsiva y, si $q_1q_2<0$, atractiva.
\end{definicion}

% Figura: ley de Coulomb (geometría) y principio de superposición.
\begin{figure}[ht]
\centering
\tdplotsetmaincoords{72}{125}
\begin{minipage}[t]{0.49\textwidth}\centering
\begin{tikzpicture}[tdplot_main_coords, scale=1.0, baseline=(current bounding box.north)]
  \draw[eje] (0,0,0) -- (3.2,0,0) node[below left] {\small $x$};
  \draw[eje] (0,0,0) -- (0,3.4,0) node[right] {\small $y$};
  \draw[eje] (0,0,0) -- (0,0,3.2) node[above] {\small $z$};
  \coordinate (Q1) at (2.6,0.4,0.6);
  \coordinate (Q2) at (0.4,2.7,2.3);
  \draw[->, azulclaro!80!azulacero, line width=0.7pt] (0,0,0) -- (Q1) node[pos=0.5, below, font=\small, text=azulprof] {$\vc{r}_1$};
  \draw[->, azulclaro!80!azulacero, line width=0.7pt] (0,0,0) -- (Q2) node[pos=0.45, left, font=\small, text=azulprof] {$\vc{r}_2$};
  \draw[vect, azulprof, line width=1.2pt] (Q1) -- (Q2) node[pos=0.62, above left, font=\small, text=azulprof] {$\vc{R}_{12}$};
  \draw[vectrojo] (Q2) -- ($(Q2)+(-0.45,0.7,0.65)$) node[above, font=\small, text=acento] {$\vc{F}_{12}$};
  \draw[vectrojo] (Q1) -- ($(Q1)+(0.45,-0.7,-0.65)$) node[left, font=\small, text=acento] {$\vc{F}_{21}$};
  \node[cargapos] at (Q1) {$+$};
  \node[cargapos] at (Q2) {$+$};
  \node[font=\small, anchor=south west] at ($(Q1)+(0.1,0.1,0.1)$) {$q_1$};
  \node[font=\small, anchor=south east] at ($(Q2)+(-0.1,0,0.15)$) {$q_2$};
\end{tikzpicture}\\[2pt]{\small\scshape (a)}
\end{minipage}\hfill
\begin{minipage}[t]{0.49\textwidth}\centering
\begin{tikzpicture}[tdplot_main_coords, scale=0.95, baseline=(current bounding box.north)]
  \draw[eje] (0,0,0) -- (3.0,0,0) node[below left] {\small $x$};
  \draw[eje] (0,0,0) -- (0,3.2,0) node[right] {\small $y$};
  \draw[eje] (0,0,0) -- (0,0,3.0) node[above] {\small $z$};
  \coordinate (P) at (1.3,1.6,1.1);
  \foreach \n/\a/\b/\c/\l in {1/2.9/0.0/0.3/{q_1}, 2/2.5/1.5/2.4/{q_2}, 3/0.2/3.1/1.8/{q_3}}{
     \coordinate (S\n) at (\a,\b,\c);
     \draw[vect, azulclaro!90!black, line width=0.7pt] (S\n) -- (P);
     \node[cargapos] at (S\n) {$+$};
     \node[font=\small, anchor=south east] at ($(S\n)+(-0.1,0,0.1)$) {$\l$};}
  \draw[vectrojo, line width=1.3pt] (P) -- ($(P)+(-0.9,0.6,0.45)$) node[above left, font=\small, text=acento] {$\vc{F}_{q}$};
  \node[cargapru] at (P) {$q$};
  \node[font=\scriptsize, text=azulprof, anchor=north] at ($(S1)!0.5!(P)$) {$\vc{R}_1$};
  \node[font=\scriptsize, text=azulprof, anchor=east] at ($(S2)!0.5!(P)$) {$\vc{R}_2$};
  \node[font=\scriptsize, text=azulprof, anchor=west] at ($(S3)!0.45!(P)$) {$\vc{R}_3$};
\end{tikzpicture}\\[2pt]{\small\scshape (b)}
\end{minipage}
\caption{(a)~Ley de Coulomb para dos cargas puntuales del mismo signo:
$\vc{R}_{12}=\vc{r}_2-\vc{r}_1$ apunta de $q_1$ a $q_2$, y la fuerza
$\vc{F}_{12}$ sobre $q_2$ es paralela a $\vc{R}_{12}$ (repulsión); la fuerza sobre $q_1$ es la
opuesta, $\vc{F}_{21}=-\vc{F}_{12}$. (b)~Superposición: la fuerza sobre la carga $q$ es la suma
vectorial de las ejercidas por cada $q_i$, con $\vc{R}_i$ el vector que une $q_i$ con $q$.}
\label{fig:coulomb}
\end{figure}


\begin{observacion}[Sobre la ley de Coulomb]\label{obs:coulomb}
\begin{enumerate}[label=(\alph*)]
  \item \emph{Acción y reacción.} Al intercambiar las cargas, $\vc{R}_{12}=-\vc{R}_{21}$ y por tanto
        $\vc{F}_{21}=-\vc{F}_{12}$: se cumple la tercera ley de Newton.
  \item \emph{La constante.} En el SI, $K=8{,}988\cdot10^{9}\,\mathrm{N\,m^{2}/C^{2}}$ y la unidad de carga,
        el culombio, se define a partir de la unidad de corriente (amperio): es la carga que atraviesa en
        un segundo la sección de un conductor por el que circula $1\,\mathrm{A}$. Se escribe habitualmente
        \[
          K=\frac{1}{4\pi\eo},\qquad
          \eo=\frac{1}{4\pi K}=8{,}85\cdot10^{-12}\,\frac{\mathrm{C^{2}}}{\mathrm{N\,m^{2}}}\;\Bigl(=\frac{\mathrm{F}}{\mathrm{m}}\Bigr),
        \]
        la \emph{permitividad del vacío}, muy próxima a la del aire seco, que es el medio donde Coulomb
        experimentó. Más adelante, en las ecuaciones de Maxwell, aparecerá la relación
        $c^{2}=1/(\eo\mu_0)$ con $c\simeq2{,}998\cdot10^{8}\,\mathrm{m/s}$ la velocidad de la luz y
        $\mu_0\approx4\pi\cdot10^{-7}\,\mathrm{H/m}$ la permeabilidad magnética del vacío~\cite{CODATA}.
  \item \emph{Alcance de validez.} La ley es experimental y su exponente está comprobado con enorme precisión
        (\cref{sec:coulomb-exp}). Aunque se estableció en condiciones macroscópicas, se aplica hasta distancias
        de $10^{-15}\,\mathrm{m}$.
  \item \emph{Tres restricciones.} (i)~\emph{Cargas puntuales}: es una idealización; las cargas deben ser pequeñas
        frente a su separación. (ii)~\emph{Cargas en reposo}: si se mueven aparece además una interacción
        magnética; el estudio de las cargas en reposo es la \emph{electrostática}. Para que se mantengan en
        reposo han de aparecer fuerzas mecánicas tales que $\sum\vF=\vc{0}$. (iii)~\emph{Cargas en el vacío}: en
        presencia de materiales aparecerían fuerzas adicionales (en los conductores, por redistribución de
        sus cargas por influencia; en los dieléctricos, por su polarización).
\end{enumerate}
\end{observacion}

% -------------------------------------------------------------------------
\subsection{Principio de superposición}\label{sec:superposicion}
% -------------------------------------------------------------------------
La fuerza entre dos cargas no cambia por la presencia de una tercera. Este hecho, resultado de la observación
experimental, es la base del \emph{principio de superposición}.

\begin{definicion}[Principio de superposición]\label{def:superposicion}
La fuerza que un sistema de cargas puntuales $q_1,\dots,q_N$ ejerce sobre una carga puntual $q$ es la suma
vectorial de las fuerzas que cada una de ellas ejerce por separado:
\begin{equation}\label{eq:superposicion}
  \vc{F}_q=\sum_{i=1}^{N}\vc{F}_{iq}=\frac{q}{4\pi\eo}\sum_{i=1}^{N}\frac{q_i}{R_i^{3}}\,\vc{R}_i
  =\frac{q}{4\pi\eo}\sum_{i=1}^{N}\frac{q_i}{R_i^{2}}\,\uh{R_i},
\end{equation}
donde $\vc{R}_i$ es el vector que une cada carga $q_i$ con la carga $q$ (\cref{fig:coulomb}b).
\end{definicion}

\begin{observacion}[Validez]\label{obs:superposicion}
La superposición se ha comprobado válida para campos tan intensos como $10^{21}\,\mathrm{V/m}$ y distancias
tan pequeñas como $10^{-10}\,\mathrm{m}$. Expresa la \emph{linealidad} del medio: vale en el vacío, pero
no necesariamente en todos los materiales. Es, junto con la ley de Coulomb, el segundo dato de partida que usaremos:
con \cref{eq:coulomb,eq:superposicion} basta para deducir el resto de la electrostática en el vacío.
\end{observacion}
